% 04_geo.tex — GEO 分析

\section{GEO（地球同步轨道）深度分析}

\subsection{辐射环境：混合型威胁}

GEO（约 35,786 km）位于范艾伦外带边缘。捕获质子几乎为零，辐射威胁主要来自：(1) 相对论性电子（$>$2 MeV）、(2) 太阳质子事件（SPE）、(3) 银河宇宙线（GCR）。

\begin{table}[H]
\centering
\caption{GEO 辐射环境关键参数}
\label{tab:geo-rad}
\resizebox{\textwidth}{!}{%
\begin{tabular}{@{}llll@{}}
\toprule
\textbf{参数} & \textbf{GEO 值} & \textbf{vs LEO} & \textbf{来源} \\
\midrule
TID (4mm Al) & 3--10~\kradyr & \textbf{0.3--0.5x 更优} & Zheng et al. (2023) \cite{zheng_scad_2023} \\
SEU 率 & $10^{-5}$--$10^{-3}$ 次/器件$\cdot$天 & 可比至 $\sim$10x 更劣 & Hands et al. (2018) \cite{hands_radiation_2018} \\
主要粒子类型 & 电子+SPE+GCR & 与 LEO 质子主导不同 & Laser2COTS \\
内带电风险 & \textbf{高}（$>$2 MeV 电子穿透屏蔽）& 更劣 & Horne et al. (2025) \cite{horne_internal_charging_2025} \\
表面充电风险 & 高（磁暴/亚暴 keV 电子）& 更劣 & Zheng et al. (2023) \cite{zheng_scad_2023} \\
GCR 通量 & $\sim$6x LEO & 更劣 & Laser2COTS \\
\bottomrule
\end{tabular}%
}
\end{table}

值得注意的是，GEO 的 TID（3--10~\kradyr）\textbf{低于} LEO（10--30~\kradyr），因为 LEO 穿越 SAA 的捕获质子贡献了大量 TID，而 GEO 无此来源。但 SEE 风险因 GCR 全通量（$\sim$6x LEO）和 SPE 直接暴露而升高。

\textbf{GEO 特有威胁}：
\begin{itemize}
    \item \textbf{内带电}：$>$2 MeV 电子穿透屏蔽介质，在介质内部积累电荷，达到击穿阈值时瞬间放电（介质击穿）。需 $\geq$ 2.8 mm Al 屏蔽控制 \cite{horne_internal_charging_2025}。
    \item \textbf{表面充电}：磁暴/亚暴期间 keV 级热等离子体注入。SCAD 卫星实测最大 -800 V 充电电压 \cite{zheng_scad_2023}。
\end{itemize}

\subsection{热环境：近乎理想}

GEO 的热环境是所有地球轨道中\textbf{最优}的：
\begin{itemize}
    \item 地球红外（$\sim$5--6~\wpm）和反照（$<$2~\wpm）可忽略
    \item 食仅在春/秋分季节发生，全年 $<$2\% 时间在食中
    \item 热循环几乎为零（无频繁进出阴影的热冲击）
    \item 350 K 散热器净散热能力约 1,040~\wpm——比 LEO 高 $\sim$2.2x
\end{itemize}

这意味着 GEO 算力卫星的散热系统设计可大幅简化：无需定向冷面散热器，无需应对频繁热循环带来的材料疲劳。

\subsection{在轨维护：GEO 的独特战略优势}

GEO 是所有高轨道中唯一具备"可维护"前景的轨道：

\begin{itemize}
    \item Northrop Grumman MEV-1 (2020) 和 MEV-2 (2021) 已成功对接 Intelsat 卫星并延寿 5 年以上
    \item DARPA RSGS 机器人有效载荷已完成 TVAC 测试（2024），即将集成至 MRV 航天器
    \item 在轨服务市场预计从 2024 年 \$2.7B 增长至 2034 年 \$8B（CAGR $\sim$11--12\%）
    \item GEO 通信卫星的 15 年设计寿命已获产业验证（Maxar 1300、Boeing 702、LM 2100 平台）
\end{itemize}

\textbf{对太空算力的关键意义}：GEO 算力卫星可定期延寿/升级。硬件寿命从 5 年延至 15--20 年 $\rightarrow$ 折旧成本下降 3--4x。这是 GEO 在综合成本评估中超越 LEO 的核心原因之一。

\subsection{太空算力部署评估}

\textbf{GEO 是综合排名第一的非 LEO 太空算力轨道}（综合成本指数 0.83，LEO 为 1.00）。主要优势：

\begin{enumerate}
    \item \textbf{15 年已验证寿命} + \textbf{在轨服务即将可用} $\rightarrow$ 全生命周期成本最优
    \item \textbf{固定覆盖 1/3 地球表面} $\rightarrow$ 区域实时推理广播、军事/情报持续监视
    \item \textbf{热设计简化} $\rightarrow$ 光伏面积需求比 LEO 少约 40\%
    \item \textbf{最可能商业化场景}：通信卫星"搭载"AI 推理载荷，对下行数据预处理压缩
\end{enumerate}

主要约束：RTT $\sim$480--600 ms 排除交互式推理（但本地处理可规避）；硬件抗辐射溢价仍是第一成本驱动因子。\textbf{关键先行事件}：DARPA RSGS 首次 GEO 在轨服务成功（预计 2027--2028）。
